\section*{الفصل \ref{ch:b1:titration-methods} --- طرائق المعايرة ومنحنياتها}

\begin{solution}{exo:b1:titration-methods:1}
\ce{H3O+ + OH- -> 2H2O}، $K = 1/K_e = 10^{14.00}$. $C = 12.4 \times
0.100/20.0 = \qty{0.0620}{mol/L}$.
\end{solution}

\begin{solution}{exo:b1:titration-methods:2}
عند \omterm{def:b1:titration-methods:titration}{التكافؤ} يكون المحلول محلول كلوريد الأمونيوم بتركيز
\qty{0.050}{mol/L}: $\mathrm{pH} = \frac12(9.25 - \log 0.050) = 5.28$.
يحويه مجال أحمر الميثيل (3.8--5.8)؛ أما برتقالي الميثيل، الذي لا يُبلغ مجاله (2.5--4.5) إلا بعد \omterm{def:b1:titration-methods:titration}{التكافؤ} مع هبوط pH، فسيتغير لونه متأخرًا.
\end{solution}

\begin{solution}{exo:b1:titration-methods:3}
\ce{MnO4- + 5Fe^2+ + 8H+ -> Mn^2+ + 5Fe^3+ + 4H2O}؛
$C = 5 \times 0.0200 \times 12.6/10.0 = \qty{0.126}{mol/L}$.
\end{solution}

\begin{solution}{exo:b1:titration-methods:4}
$C = 0.0100 \times 14.2/50.0 = \qty{2.84e-3}{mol/L}$ من الكالسيوم
والمغنيسيوم معًا. يتفاعل الاثنان قبل \omterm{def:b1:titration-methods:titration}{نقطة النهاية}، فلا تعطي إلا مجموعهما:
إنها \omterm{def:b1:titration-methods:kinds}{معايرة متزامنة}.
\end{solution}

\begin{solution}{exo:b1:titration-methods:5}
0~mL: \omterm{def:b1:predominant-reaction:strong-acid}{حمض ضعيف}، $\mathrm{pH} = \frac12(4.76 + 1.00) = 2.88$. 5.0~mL:
نصف \omterm{def:b1:titration-methods:titration}{التكافؤ}، 4.76. 10.0~mL: 8.73
(\cref{prop:b1:titration-methods:ph-at-equivalence}). 15.0~mL: هيدروكسيد
فائض \qty{0.50}{mmol} في \qty{25.0}{mL}، $[\ce{OH-}] = 0.020$،
$\mathrm{pH} = 12.30$. وكلها تتفق مع المنحنى المحسوب في حدود 0.01.
\end{solution}

\begin{solution}{exo:b1:titration-methods:6}
قيمتا $\mathrm{p}K_a$ 2.15 و7.21 تختلفان بأكثر من 4، وكذلك 7.21
و12.34: تكافؤان، عند 10.0 و\qty{20.0}{mL}. الأول: محلول
\ce{H2PO4-}، $\mathrm{pH} \approx \frac12(2.15 + 7.21) = 4.68$؛ والثاني:
\ce{HPO4^2-}، $\frac12(7.21 + 12.34) = 9.78$ (يعطي المنحنى الدقيق 4.71
و9.66، إذ تهمل الصيغتان التخفيف والماء). أما الحمضية الثالثة
($\mathrm{p}K_a$ 12.34) فقريبة جدًا من حمضية الماء: فلا يتفاعل الهيدروكسيد المضاف
إلا جزئيًا ويرتفع pH دون قفزة.
\end{solution}

\begin{solution}{exo:b1:titration-methods:7}
حمض الكلوريدريك: $6.2 \times 0.100/10.0 = \qty{0.062}{mol/L}$؛ وحمض
الإيثانويك: $(14.6 - 6.2) \times 0.100/10.0 = \qty{0.084}{mol/L}$. وفي المنتصف يكون نصف
حمض الإيثانويك قد عُوير: $\mathrm{pH} = \mathrm{p}K_a = 4.76$.
\end{solution}

\begin{solution}{exo:b1:titration-methods:8}
$V_{\mathrm{eq}} = \qty{10.0}{mL}$. 2.0~mL: $E = 0.77 + 0.059\log(2/8) =
0.73$~V؛ 9.0~mL: $0.77 + 0.059\log 9 = 0.83$~V؛ 11.0~mL: برمنغنات
فائضة \qty{0.020}{mmol}، و\ce{Mn^2+} \qty{0.200}{mmol}، $E = 1.51 +
\frac{0.059}{5}\log 0.10 = 1.50$~V؛ 20.0~mL: 1.51~V.
\end{solution}

\begin{solution}{exo:b1:titration-methods:9}
0~mL: $\kappa = (349.8 + 76.2) \times 0.0100 = \qty{4.26}{mS/cm}$.
10.0~mL: \ce{Na+} و\ce{Cl-} بتركيز $1.00/110 = \qty{9.09e-3}{mol/L}$،
$\kappa = (50.3 + 76.2) \times \num{9.09e-3} = 1.15$. 20.0~mL: \ce{Na+}
$0.0167$، و\ce{Cl-} و\ce{OH-} $0.00833$~mol/L، $\kappa = 0.84 + 0.635 + 1.65
= 3.13$~mS/cm. الميلان (مع إهمال التخفيف، \qty{1}{mL} يضيف
\qty{1.0e-3}{mol/L}): $(50.3 - 349.8) \times 10^{-3} = -0.30$
و$(50.3 + 198.3) \times 10^{-3} = +0.25$~mS/cm لكل mL.
\end{solution}

\begin{solution}{exo:b1:titration-methods:10}
عند $V_{\mathrm{eq}} - v$ يكون الحمض الفائض $C_Bv$، وعند $V_{\mathrm{eq}} + v$
تكون القاعدة الفائضة $C_Bv$؛ وفي الحجم نفسه، يساوي $[\ce{H3O+}]$ قبل \omterm{def:b1:titration-methods:titration}{التكافؤ}
$[\ce{OH-}]$ بعده، فيكون $\mathrm{pH}(V_{\mathrm{eq}} - v) +
\mathrm{pH}(V_{\mathrm{eq}} + v) = 14$: تناظر حول النقطة (\babelsublr{V$_{\mathrm{eq}}$}،
7). القفزة: \qty{1.0e-5}{mol} فائضة في \qty{20}{mL}، من pH 3.30 إلى 10.70، أي 7.4
وحدة. وبالقسمة على 100: $\num{1.0e-7}$~mol في \qty{20}{mL}، من pH 5.3 إلى 8.7،
أي 3.4 وحدة: \omterm{def:b1:titration-methods:titration}{نقطة نهاية} أقل حدة بكثير.
\end{solution}

\begin{solution}{exo:b1:titration-methods:11}
قبل \omterm{def:b1:titration-methods:titration}{التكافؤ}: \ce{Fe^3+} المتكوّن $= 5C_{\mathrm{Mn}}V$، و\ce{Fe^2+} الباقي $=
5C_{\mathrm{Mn}}(V_{\mathrm{eq}} - V)$، فيكون $E = 0.77 + 0.059\log\frac{V}{V_{\mathrm{eq}}
- V}$: وعند $V_{\mathrm{eq}}/2$، $E = 0.77$. وبعده: فائض \ce{MnO4-} $\propto
V - V_{\mathrm{eq}}$، و\ce{Mn^2+} $\propto V_{\mathrm{eq}}$: $E = 1.51 +
\frac{0.059}{5}\log\frac{V - V_{\mathrm{eq}}}{V_{\mathrm{eq}}}$ (pH 0)؛ وعند
$2V_{\mathrm{eq}}$، 1.51~V. وعند \omterm{def:b1:titration-methods:titration}{التكافؤ} $(0.77 + 5 \times 1.51)/6 =
1.39$~V. وعند pH 1 تنخفض مزدوجة البرمنغنات بمقدار $0.059 \times 8/5 =
0.094$~V: $E_{\mathrm{eq}} = (0.77 + 5 \times 1.416)/6 = 1.31$~V.
\end{solution}

\begin{solution}{exo:b1:titration-methods:12}
للتفاعل \ce{NH4+ + OH- -> NH3 + H2O} الثابت $K = 10^{14.00 - 9.25} = 10^{4.75}$: فلا يكتمل
التفاعل قرب \omterm{def:b1:titration-methods:titration}{التكافؤ} والقفزة صغيرة،
نحو pH 11 (pH \omterm{def:b1:titration-methods:titration}{التكافؤ} $14 - \frac12(4.75 + 1.30) = 11.0$). \omterm{def:b1:titration-methods:kinds}{المعايرة
العكسية}: أضف $n_1$ من هيدروكسيد الصوديوم (بفائض)، واغلِ لطرد
الأمونيا المتكوّنة، ثم عاير الهيدروكسيد الباقي بحمض الكلوريدريك
($n_2$، قفزة حادة قوي--قوي): $n(\ce{NH4+}) = n_1 - n_2$.
\end{solution}

\begin{solution}{pb:b1:titration-methods:1}
\textbf{1.} \ce{C6H6O6 + 2H+ + 2e- -> C6H8O6}.
\textbf{2.} \ce{C6H8O6 + I2 -> C6H6O6 + 2I- + 2H+}.
\textbf{3.} 0 و\babelsublr{$-$I}.
\textbf{4.} \ce{I2 + 2S2O3^2- -> 2I- + S4O6^2-}، $\log K = 2(0.53 - 0.02)/0.059
= 17$.
\textbf{5.} يحتاج اللون الأزرق إلى اليود؛ فهو يختفي مع آخر
يود، عند \omterm{def:b1:titration-methods:titration}{التكافؤ}.
\textbf{6.} الفقد بالهواء كان سيخفض النتيجة؛ أما إذا أضيف اليود دفعة واحدة
وبفائض، فإنه يؤكسد حمض الأسكوربيك بسرعة وبشكل تام.
\textbf{7.} تفاعلًا سريعًا \omterm{def:b1:extent-q-and-k:quantitative}{وكمّيًا} طوال \omterm{def:b1:titration-methods:titration}{المعايرة}،
دون أكسدة بالهواء أثناء إضافة \omterm{def:b1:titration-methods:titration}{المحلول المعايِر} ببطء، \omterm{def:b1:titration-methods:titration}{ونقطة نهاية}
عند أول فائض من اليود.
\textbf{8.} اليود الفائض $n_R$؛ والتفاعل مع حمض الأسكوربيك؛ \omterm{def:b1:titration-methods:titration}{ومعايرة}
اليود الباقي بالثيوكبريتات.
\textbf{9.} وإلا بقي بعض حمض الأسكوربيك ولم يكن فائض
اليود، المعدوم، ليدل على مقداره. وهنا
تفاعل \qty{2.775e-4}{mol} من أصل \qty{5.000e-4}{mol}: فهو فائض فعلًا.
\textbf{10.} تتغير محاليل الثيوكبريتات ببطء مع الوقت؛ وتركيزها
يدخل مباشرة في النتيجة.
\textbf{11.} 10.
\textbf{12.} ماصة معيارية سعتها \qty{20.00}{mL}.
\textbf{13.} $20.00 \times 10^{-3} \times 0.0250 = \qty{5.000e-4}{mol}$.
\textbf{14.} $8.90 \times 10^{-3} \times 0.0500 = \qty{4.450e-4}{mol}$.
\textbf{15.} النصف: \qty{2.225e-4}{mol}.
\textbf{16.} $5.000 - 2.225 = \qty{2.775e-4}{mol}$.
\textbf{17.} واحد لواحد: \qty{2.775e-4}{mol}.
\textbf{18.} $\times 10$: \qty{2.775e-3}{mol}، $\times 176.0$:
\qty{0.488}{g}.
\textbf{19.} أدنى من اللصاقة بنحو \qty{2}{\percent}.
\textbf{20.} $n = \bigl(C_IV_I - \frac12C_TV_T\bigr)\,V_{\mathrm{flask}}/V_{\mathrm{aliquot}}$.
\textbf{21.} اليود المُدخل: الماصة $0.03/20.00 = \qty{0.15}{\percent}$
والتركيز \qty{0.2}{\percent}، أي نحو \qty{0.25}{\percent}، أو
\qty{1.3e-6}{mol}. والفائض: السحاحة $0.05/8.90 = \qty{0.56}{\percent}$
والتركيز \qty{0.2}{\percent}، أي نحو \qty{0.60}{\percent}، أو
\qty{1.3e-6}{mol}.
\textbf{22.} تتجمع الارتيابات المطلقة، بينما الفرق
أصغر من كل من الكميتين: $\sqrt{1.3^2 + 1.3^2} \times 10^{-6} =
\qty{1.8e-6}{mol}$، أي \qty{0.66}{\percent} من \qty{2.775e-4}{mol}، وهو أكثر من
كل من الحدين.
\textbf{23.} مع الحوجلة (\qty{0.1}{\percent}) وماصة الجزء المأخوذ
(\qty{0.2}{\percent}): $\sqrt{0.66^2 + 0.1^2 + 0.2^2} = \qty{0.70}{\percent}$.
\textbf{24.} $0.488 \pm 0.003$~g: \textbf{\qty{0.49}{g} من حمض الأسكوربيك
في كل قرص}.
\end{solution}
